% --- INICIO DEL MÓDULO: 04_conclusiones.tex ---

\section{Conclusiones}

El incremento de \(Sps\) mejora la representación temporal de la señal y permite observar con mayor claridad la estructura de la PSD, sin modificar el ancho de banda fundamental determinado por la tasa de bits.


Las señales obtenidas a partir de \texttt{rana.jpg} y \texttt{sonido.wav} presentan distribuciones espectrales distintas a las de una fuente binaria aleatoria ideal, debido a la correlación estadística propia de cada archivo.


Las modificaciones realizadas sobre el filtro interpolador mostraron que la codificación de línea y la conformación de pulso alteran la componente DC y la ocupación espectral de la señal, por lo que la elección del pulso influye directamente en la eficiencia de transmisión.
